\subsection{A concrete polarized theta germ at the height-three prime}

\begin{proposition}[DP49: actual Jacobian theta germs and nonzero primitive paths]\label{DP49}
Let $R=W(\mathbf F_{11^3})$, $K=\operatorname{Frac}R$, and let
$J=\operatorname{Jac}(y^2=x^7-1)$ with DP22's good integral model.
Let $a$ be the covariant Jacobian automorphism induced by
$(x,y)\mapsto(\zeta_7x,y)$, and let $\lambda$ be its canonical principal polarization.
Put
\[
 w_0=-(a+a^2+a^4).
\]
Then $a^\dagger=a^{-1}$, $w_0+w_0^\dagger=1$, and $w_0w_0^\dagger=2$.
The height-three factor $G$ on the orbit $(3,5,6)$ has dimension one; the other factor $H$ has dimension two and is integrally its actual Cartier dual.
Their integral displays are the arbitrary-unit singleton/complement models of DP46.
The polarized formal completion is the actual hyperbolic block $G\times H$ of DP47.

Let $\mathcal P_\lambda$ be the symmetric polarization biextension and
$\mathcal L_{w_0}=(1,w_0)^*\mathcal P_\lambda$.
This trace-one line has associated biextension $\mathcal P_\lambda$.
Fix an integer $\kappa>0$ prime to $11$ and a nonzero integral section germ $f$ of
$\mathcal L_{w_0}^{\otimes\kappa}$ in an integral frame.
Use the unit-arrow logarithm coordinates $X,Y_1,Y_2$.
The corresponding tangent directions have original Jacobian CM characters $3,2,1$.
For a logarithmic line
\[
 X=ut,\qquad (Y_1,Y_2)=(v_1t,v_2t),
 \qquad u\in R^\times,\quad (v_1,v_2)\ \hbox{primitive},
\]
suppose the integral algebraic prefactor has nonzero restriction.
In the invariant product canonical gauge, writing the normalized germ as
$\sum_n A_nt^n$, one has
\[
 v_{11}(n!A_n)\ge c_kn-\frac{s_{11}(n)}{10},
 \qquad
 \liminf_{n\to\infty}\frac{v_{11}(n!A_n)}n=c_k,
\]
where $s_{11}(n)$ is the base-$11$ digit sum and
\[
\begin{array}{c|c|c|c}
\text{dual direction}&k&w_k&c_k=1/10-w_k\\ \hline
v_1\in R^\times&2&121/1220&1/1220\\
v_1\in11R,\ v_2\in R^\times&1&11/120&1/120 .
\end{array}
\]
The integer $k$ is the reduced dual tangent vector's restricted-$11$ nilpotency length.

For any rational canonical gauge let $ct^2$ be its remaining quadratic correction after the integral-frame adjustment of DP47.
With $r_c=(1/10-v_{11}(c))/2$, or $r_0=-\infty$, the exact coefficient threshold is
$R_c=\max(w_k,r_c)$ and the exact factorial-slope limit is $1/10-R_c$.
An integral $c$ leaves the displayed slope unchanged; if $v_{11}(c)=-M<0$, the limit is
$(1/10-M)/2<0$.

Finally, for every fixed $\bar u\in(R/11R)^\times$ and every nonzero
$\bar v\in(R/11R)^2$, some lifts of these exact residues give a nonzero restriction of $f$.
Thus both displayed rates are attained in the product gauge for a nonzero germ.
This is existence in every primitive residue class, not nonvanishing of a path prescribed in advance.
\end{proposition}

\begin{proof}
DP22 supplies the actual smooth genus-three curve and Jacobian abelian scheme, with the integral cyclotomic action.
The differentials $dx/y,x\,dx/y,x^2dx/y$ have characters $1,2,3$.
The covariant Abel--Jacobi convention
$a[P-\infty]=[\sigma(P)-\infty]$ gives the same scalar characters on the tangent space: ordinary transposition retains each scalar.
The curve automorphism fixes infinity and preserves effective divisors, so it preserves the canonical theta polarization.
Thus $a^\vee\lambda a=\lambda$, giving $a^\dagger=a^{-1}$.

Since $11\equiv4\pmod7$ has order three, the integral unramified cyclotomic algebra has the two conjugate factors
$\{1,2,4\}$ and $\{3,5,6\}$.
Their heights and dimensions are those of DP22.
A nonzero etale quotient of either height-three factor would have a rational Tate space over the unramified cubic field, hence positive height divisible by three.
It would exhaust the factor, contradicting its positive dimension.
Both factors are therefore connected; the same applies to their duals.
These are the functorial connected--etale and dimension statements of
\cite[\S2.2, (4); \S2.3, Proposition 3]{Tate67}.

Here the integral unit-arrow normal form is also explicit.
Over $R_n=R/11^n$, use the actual display from
\cite[Theorems 5 and 9, pp.~131--132]{Zink}.
Its rank-six Witt module splits by
\[
 e_b=\frac17\sum_{r=0}^6[\zeta_7]^{-br}a^r,\qquad 1\le b\le6,
\]
into free eigenspaces; the zeroth summand is zero by $\Phi_7(a)=0$.
Frobenius sends character $b$ to $11b$.
Its linearization becomes invertible after inverting $11$:
in a normal decomposition the isomorphism $V^{-1}\oplus F$ becomes that linearization after its $L$ columns are multiplied by $11$.
The coefficient ring after this inversion is nonzero, as it maps to
$W(\mathbf F_{11^3})[1/11]$.
Consequently ranks are equal along each orbit.
Both orbit ranks are positive because the tangent quotient contains characters $1$ and $3$.
Their common ranks $r_1,r_2$ satisfy $3r_1+3r_2=6$, so both are one.

The exact Hodge quotient $P_n/Q_n=\operatorname{Lie}J$ now gives
$Q_n\cap P_{b,n}=I_nP_{b,n}$ for $b=1,2,3$ and
$Q_n\cap P_{b,n}=P_{b,n}$ for $b=4,5,6$.
In the orbit $(3,5,6)$ there is exactly one selected line.
Choose generators in this order.
The normal-decomposition isomorphism is monomial on them, so every arrow coefficient is a Witt unit:
\[
 Fe_0=u_0e_1,\qquad Fe_1=11u_1e_2,\qquad
 Fe_2=11u_2e_0 .
\]
Generators lift compatibly through the surjective Witt reductions, and units remain units.
Their inverse limits give the actual integral unit-arrow model over $R$, exactly as in DP45; no isogeny denominator or nonunit tangent map has been introduced.
The dual display is the complementary reciprocal-unit display
\cite[Definition 19 and (22), p.~143]{Zink}.

Rosati conjugation exchanges the two integral idempotents.
For two points in the same factor, balance by its idempotent identifies the polarization biextension with its zero-axis restriction.
The diagonal restrictions are thus trivial as biextensions.
The principal polarization is perfect on the $11$-divisible group and identifies $H$ with the actual $G^D$.
Symmetry identifies the opposite cross restriction with its transpose; biadditivity gives the hyperbolic product.
This verifies DP47's integral geometric hypotheses.
The ordinary dual action on its two tangent positions retains characters $5,6$.
The original Jacobian action on $H$ is $(a_G^{-1})^D$, so those become $-5=2$ and $-6=1$ modulo seven.
The unit arrows give the nonzero restricted-power chain from character $2$ to character $1$, followed by zero.
This proves the intrinsic description of $k$.

For the trace-one calculation put $\eta=a+a^2+a^4$.
Then
\[
 \eta+\eta^\dagger=\sum_{r=1}^6a^r=-1,\qquad
 \eta\eta^\dagger=3+\sum_{r=1}^6a^r=2.
\]
Thus $w_0=-\eta$ has trace one and norm two.
Also $\eta^2+\eta+2=0$, so $\delta=1+2\eta$ satisfies
$\delta^2=-7$ and $w_0=(1-\delta)/2$.
Its polynomial expression proves integrality directly.

The associated biextension of $\mathcal L_{w_0}$ is
\[
 \mathcal P_\lambda(z,w_0z')\otimes
 \mathcal P_\lambda(z',w_0z).
\]
Symmetry and adjunction turn its second factor into
$\mathcal P_\lambda(z,w_0^\dagger z')$.
The trace-one relation and addition in the second argument give exactly
$\mathcal P_\lambda(z,z')$.
Hence the trace-one line induces the actual principal polarization.

DP47 now identifies the normalized germ, in an integral frame, with
\[
 \widetilde f(S,Y)B(S,Y)^{-\kappa}
       \exp Q_\kappa(\log_GS,\log_HY),
\]
where $\widetilde f$ is integral and $Q_\kappa$ is quadratic.
The coefficient $-\kappa$ is a unit.
On the prescribed logarithmic line the quadratic is $ct^2$.
DP46 gives the exact primitive threshold and factorial limit for the actual unit-arrow model.
DP48's proof applies without change, since it uses precisely that threshold, bounded nonzero restricted prefactors and the bounded-multiplier lemma.
It yields $R_c=\max(w_k,r_c)$ and the stated limit.
In the invariant product gauge $Q_\kappa=0$.
Substituting $p=11$ and $k=1,2$ gives all entries of the table and both factorial inequalities.

It remains to prove that the nonzero restrictions exist.
Let $f_1$ be the integral germ including its unit frame factor.
Inverse logarithms define an automorphism of $K[[X,Y_1,Y_2]]$, so
$F=f_1(E_G(X),E_H(Y))$ is nonzero.
Let $P_d$ be its first nonzero homogeneous polynomial.
The inverse logarithms have identity tangent, so this is the corresponding leading polynomial of $f_1$ as well.
The lifts of any fixed primitive residues $(\bar u,\bar v)$ form a nonempty full $11$-adic open box.
A nonzero polynomial over $K$ cannot vanish on it:
in one variable it has only finitely many roots, whereas an open ball is infinite; induction on the number of variables applies this fact to its coefficient polynomials.
Choose lifts with $P_d(u,v)\ne0$.
Then the coefficient of $t^d$ in $F(ut,vt)$ is nonzero.
The residue class fixes the same restricted-power length $k$, proving the last assertion.

For a nonzero algebraic section on the integral Jacobian model, its local germ is nonzero because vanishing on a nonempty open set would force vanishing at the generic point and everywhere.
Its formal completion remains nonzero by the local Krull intersection theorem: for a proper ideal in a Noetherian local ring, the intersection of all its powers on a finite module is zero
(Stacks Project, Lemma 10.51.4, Tag 00IP,
\url{https://stacks.math.columbia.edu/tag/00IP}).
This requires neither a nonzero value at the origin nor a primitive initial homogeneous part.
\end{proof}

DP49 supplies actual normalized theta-germ rates on a concrete higher-height polarized Jacobian block.
The best ordinary primitive mixed-line rate is $1/(11^2-1)$.
It does not identify the global construction's prescribed paths, selected harmonic polynomials or minors with these ordinary lines, and the polynomial argument does not establish nonvanishing on a previously fixed path.
The Jacobian still fails the old induced-type host condition of DP25.
The local primitive Poincar\'e pullback, however, is an actual proved identity in this example.

